\setcounter{aufgabe}{9}
\hypertarget{e2}{}\einheitenkopf{Einheit 2 von 5 · Satz vom Nullprodukt}

\begin{aufgabe}{Nullprodukt -- Steht rechts eine Null? Kreuze an. Löse nichts.}
\begin{teile}
\teil $x \cdot (x - 4) = 0$ \hfill \janein
\teil $(x + 2) \cdot (x - 9) = 12$ \hfill \janein
\teil $(x - 6) \cdot (x + 1) = 0$ \hfill \janein
\teil $x \cdot (x + 11) = 24$ \hfill \janein
\teil $5x \cdot (x - 3) = 0$ \hfill \janein
\end{teile}
\end{aufgabe}

\begin{aufgabe}{Nullprodukt -- rechnen. Setze jeden Faktor einzeln null und gib beide
Lösungen an.}
\beispiel{(x - 3) \cdot (x + 8) &= 0 \\
x - 3 &= 0 && \text{oder}\quad x + 8 = 0 \\
x_1 &= 3 && x_2 = -8}
\begin{gleichungsraster}[2]
\gl{(x - 2) \cdot (x - 9) = 0} & \gl{(x - 4) \cdot (x - 11) = 0} \\
\gl{(x - 1) \cdot (x - 6) = 0} & \gl{(x - 5) \cdot (x - 12) = 0} \\
\gl{(x + 7) \cdot (x - 3) = 0} & \gl{x \cdot (x + 4) = 0} \\
\gl{(x - 6) \cdot (x - 6) = 0} & \\
\end{gleichungsraster}
Hier steht kein fertiges Produkt. Klammere zuerst $x$ aus.
\begin{gleichungsraster}[3]
\gl{x^2 + 9x = 0} & \gl{x^2 - 13x = 0} \\
\gl{4x^2 + 20x = 0} & \gl{3x^2 - 12x = 0} \\
\end{gleichungsraster}
\hfill (P10 2020)
\end{aufgabe}

\begin{aufgabe}{Nullprodukt -- umformen. Rechts steht keine Null, der Satz gilt hier
nicht. Multipliziere aus, bring alles auf eine Seite und ordne nach $x^2$, $x$ und
Zahl. Lösen kommt in Einheit 3.}
\begin{gleichungsraster}[3]
\gl{x \cdot (x + 6) = 27} & \gl{(x - 2) \cdot (x + 5) = 18} \\
\gl{x \cdot (x - 9) = -14} & \\
\end{gleichungsraster}
\end{aufgabe}

\begin{aufgabe}{Nullprodukt -- Probe. Setze die Zahl für $x$ ein, rechne zuerst die
Klammer aus und kreuze an.}
\begin{teile}
\teil $x \cdot (x + 10) = -24$ \quad für $x = -6$ \hfill \kreuz{wahr}\kreuz{falsch}
\teil $(x - 7) \cdot (x + 2) = -20$ \quad für $x = 3$ \hfill \kreuz{wahr}\kreuz{falsch}
\teil $(x + 5) \cdot (x - 1) = 9$ \quad für $x = -2$ \hfill \kreuz{wahr}\kreuz{falsch}
\end{teile}
\end{aufgabe}

\begin{aufgabe}{Nullprodukt -- Fehler finden. Sara löst so:}
\rechnung{x^2 &= 11x &&\mid : x \\ x &= 11}
Finde den Fehler, benenne ihn und korrigiere die Rechnung.

\begin{teile}
\teil Löse $x^2 = 6x$.
\end{teile}
\end{aufgabe}

\begin{aufgabe}{Nullprodukt -- Begründen. Antworte ohne zu rechnen.}
\begin{teile}
\teil Begründe, warum ein Produkt null ist, sobald einer der beiden Faktoren null ist.
\teil Lena will $x \cdot (x - 8) = 0$ lösen und teilt dafür beide Seiten durch $x$.
Erkläre, warum sie dabei eine Lösung verliert.
\teil $(x + 3) \cdot (x - 3) = 0$ und $(x + 3) \cdot (x + 3) = 0$: Welche Gleichung
hat zwei Lösungen, welche nur eine? Begründe.
\end{teile}
\end{aufgabe}

\begin{aufgabe}{Nullprodukt -- Lösung ankreuzen.}
\begin{teile}
\teil Welche Zahlen lösen $x \cdot (x - 6) = 0$? Kreuze alle an, die passen. \\
\kreuz{$x = 0$} \\ \kreuz{$x = 6$} \\ \kreuz{$x = -6$} \\ \kreuz{$x = 1$}
\teil Welche Zahl löst $x \cdot (x + 7) = -12$? Kreuze an. \hfill (P10 2025) \\
\kreuz{$x = 3$} \\ \kreuz{$x = -3$} \\ \kreuz{$x = -12$} \\ \kreuz{$x = 4$}
\end{teile}
\end{aufgabe}
